% S4 proofs for Section 5
\section{Proofs for Section 5: signs, reinforcement and monotonicity}\label{sec:proofs-signs}

Finite-population notation: $p_i=N_i/\sum_jN_j$, $\E_p$, $\mu_r=\E_p[X_r]$,
$w^r_i=p_iX_{ir}/\mu_r$, $w^{\neg r}_i=p_i(1-X_{ir})/(1-\mu_r)$, $B_s=\E_{w^s}\beta_s$.
Where $X_{ir}=1$, the absent ``rest'' group is given any bounded value; its weight is
zero. A column is fixed and dropped from the notation.

\subsection{Identities}

\begin{proofof}{Lemma~\ref{P-lem:M1}}
(a) By A1, $T_i=\sum_sX_{is}\beta_{is}$, so $\E_p[XT]=\sum_s\E_p[XX_s\beta_s]$. For each
$s$, $\E_p[XX_s\beta_s]=\mu_s\E_{w^s}[X\beta_s]=\mu_s(\E_{w^s}[X]B_s+\Cov_{w^s}(X,\beta_s))$,
and $\mu_s\E_{w^s}[X]=\E_p[XX_s]=Me_s$. Hence $\E_p[XT]=MB+\Delta$, and if $M$ is
invertible $g=B+M^{-1}\Delta$. Also
$\one^\top\Delta=\sum_s\mu_s\Cov_{w^s}(\one^\top X,\beta_s)=0$ because $\one^\top X=1$ is
constant.

(b) $\mu_r\NM_r=\E_p[X_rT]=\E_p[X_r^2\beta_r]+\sum_{s\ne r}\E_p[X_rX_s\beta_s]$, and,
using $X_r=1-\sum_{s\ne r}X_s$,
$\mu_rB_r=\E_p[X_r\beta_r]=\E_p[X_r^2\beta_r]+\sum_{s\ne r}\E_p[X_rX_s\beta_r]$.
Subtracting, $\mu_r(\NM_r-B_r)=\sum_{s\ne r}\E_p[X_rX_s(\beta_s-\beta_r)]=-\delta^r_B$.

(c) $\beta^{\neg r}_i$ is a convex combination of the $\beta_{is}\in[0,1]$, $s\ne r$, so it
lies in $[0,1]$, and $T_i=X_{ir}\beta_{ir}+(1-X_{ir})\beta^{\neg r}_i$. Its group mean is
$\E_p[(1-X_r)\beta^{\neg r}]/(1-\mu_r)=\sum_{s\ne r}\mu_sB_s/(1-\mu_r)=B_{\neg r}$. So
$(X_r,\beta_r,\beta^{\neg r})$ is a two-group population with group means $B_r$ and
$B_{\neg r}$. Its between association is
$\E_p[X_r(1-X_r)(\beta_r-\beta^{\neg r})]=\E_p[X_r((1-X_r)\beta_r-\sum_{s\ne r}X_s\beta_s)]=\delta^r_B$,
and its within association is $\delta^r_W$ by definition. EGGHM's two-group identities
(their Propositions~4 and~6), with individual-level variance $\mu_r(1-\mu_r)$ and
between-unit variance $\Var_p(X_r)$, give the two formulas.

(d) $B$ and $\delta_B$ are linear in $\beta$ with design weights.
$\Cov_w(\beta_s,X_r)=\sum_iw_i\beta_{is}(X_{ir}-\E_wX_r)$ is linear in $\beta$. In
$\delta^r_W$, $(1-\mu_r)w^{\neg r}_i\beta^{\neg r}_i=p_i\sum_{s\ne r}X_{is}\beta_{is}$, so the
division by $1-X_{ir}$ cancels.
\end{proofof}

\subsection{Sign restrictions on realised tables}

\begin{proofof}{Proposition~\ref{P-prop:N1}}
A table of non-negative real masses $n_{irc}=N_iX_{ir}\beta_{irc}$ reproduces the
margins iff $\beta_i\in S(X_i,T_i)$, and it satisfies the restrictions iff $\beta\in F$
(Lemma~\ref{P-lem:M1}(d)). $F$ is a bounded polyhedron, because each $S(X_i,T_i)$ is
compact and the restrictions are linear, and $\varphi$ is affine in $\beta$. So
$\varphi(F)$ is a compact interval with both endpoints attained, which is (a); (b) is the
definition; (c) applies (a) to $F_+$ and $F_-$. (d) $\varepsilon$ enters only the
right-hand side, so the sets are nested componentwise, and right-hand-side LP duality
gives continuous piecewise-linear endpoints where $F(\varepsilon)\ne\emptyset$. A
conclusion such as ``upper endpoint $\le c$'' holds on a sublevel set of such a
function, which need not be a product set.
\end{proofof}

\begin{proofof}{Proposition~\ref{P-prop:integer}}
Let $Y_1$ be the number of group-1 members with the outcome and $Y$ the total number
with the outcome. Then $D=Y_1/N_1-(Y-Y_1)/N_0=Y_1(1/N_1+1/N_0)-Y/N_0$. Each unit's
group-1 count with the outcome ranges over an integer interval, so their sum takes
every integer between its extremes, and those extremes give the Duncan--Davis
extremes. Hence the integer gap set is the lattice $\{-Y/N_0+j(1/N_1+1/N_0)\}$ inside
$[D^-,D^+]$, with both endpoints attained, and $1/N_1+1/N_0=1/(N\mu(1-\mu))$. Each sign
row of Proposition~\ref{P-prop:tworowsign} is a half-line in $D$, so the integer set
under sign rows is the continuous set intersected with the lattice; the endpoints
differ by less than one step, and integer emptiness can occur only when the continuous
set is shorter than one step.
\end{proofof}

\emph{The four-person example.} Three units of size 4 with group-1 counts 1, 2, 3 and
outcome counts 0, 1, 0. Each group has six members, so the integer gaps are exactly
$\{-1/6,1/6\}$. By symmetry $D_{\ER}=D_{\NM}=0$, so the continuous CR set is $\{0\}$ and
no integer table satisfies CR.

\begin{proofof}{Proposition~\ref{P-prop:tworowsign}}
By Lemma~\ref{P-lem:M1}(c) with $R=2$, $\delta_W=(D_{\ER}-D)\Var_p(X)$ and
$\delta_B=(D-D_{\NM})\mu(1-\mu)$, which gives (a) because $\Var_p(X)>0$. $D$ is affine in
$\beta$ over the product of the unit polytopes, so its range is $[D^-,D^+]$ with every
value attained (Proposition~\ref{P-prop:N1}(a)). (b) $F=\{D\le D_{\ER}\}$, whose image is
$[D^-,D^+]\cap(-\infty,D_{\ER}]$, empty iff $D_{\ER}<D^-$. (c) Put
$\beta_{i1}=\beta_{i2}=T_i$ in every unit. This lies in $S(X_i,T_i)$ (and in the unit
restriction set, because that is the whole simplex by hypothesis), gives $\delta_B=0$,
and hence $D=D_{\NM}$.

(d) $D_{\NM}=\gamma D_{\ER}$ is EGGHM's Lemma~1; directly,
$D_{\NM}=\Cov_p(X,T)/(\mu(1-\mu))$ and $D_{\ER}=\Cov_p(X,T)/\Var_p(X)$. $\gamma>0$ since
$\Var_p(X)>0$, and $\gamma\le1$ because $\mu(1-\mu)-\Var_p(X)=\E_p[X(1-X)]\ge0$, with
equality iff every unit with $p_i>0$ is pure. Suppose $D_{\ER}>0$. The negative branch
requires $D\ge D_{\ER}$ and $D\le D_{\NM}=\gamma D_{\ER}\le D_{\ER}$, so
$D=D_{\ER}=D_{\NM}$, which forces $\gamma=1$; so it is empty when $\gamma<1$. When
$\gamma=1$, every positive-weight unit is pure, each unit's rate for its only group
equals $T_i$, $D$ is point-identified and equals $D_{\ER}=D_{\NM}$, and both branches are
non-empty with gap image $\{D_{\ER}\}$. The positive branch is
$\{D_{\NM}\le D\le D_{\ER}\}\cap[D^-,D^+]=[D_{\NM},\min\{D^+,D_{\ER}\}]$, non-empty because
$D^-\le D_{\NM}\le D_{\ER}$ and $D_{\NM}\le D^+$; on it $D\ge D_{\NM}>0$. The case
$D_{\ER}<0$ is symmetric, and $D_{\ER}=0$ gives $D_{\NM}=0=D$ on both branches.
\end{proofof}

\emph{The equal-rates witness in R$\times$C.} Setting $\beta_{ir}=T_i$ for every row
makes every pairwise between association zero, so it satisfies every weak between
sign and, for each pairing, the CR union. Which branch contains it depends on its
within association. Example: equal weights on $(.2,.5,.3)$, $(.5,.2,.3)$, $(.8,.1,.1)$
with every group's rate equal to $X_1$, so $T=X_1$. For group 1 against the rest,
$\ER=1$, $\gamma=.24$, and the witness has $D=\NM=.24$, $\delta_B=0$ and
$\delta_W=\Var_p(X_1)(\ER-D)=.0456>0$: it lies in the positive branch only.

\begin{proofof}{Proposition~\ref{P-prop:joint}}
Adding constraints shrinks the feasible set. For the two-column example: two
$2\times2$ units at $x=.25$ and $.75$ with complementary columns have
$\delta_{W,2}=-\delta_{W,1}$, so the same weak within sign on both columns forces both to
zero, hence $D=D_{\ER}$; opposite signs on complementary columns duplicate one
restriction.
\end{proofof}

\begin{proofof}{Proposition~\ref{P-prop:Krows}}
(a) For each rung, $m_i=m_{\rm id}(V_i)+L_ik$ with $L_i$ fixed by the design
(Lemmas~\ref{P-lem:decomp} and~\ref{P-lem:D1}). Each association is linear in
$(m_i)_i$ with design coefficients (Lemma~\ref{P-lem:M1}(d) with $m$ in place of
$\beta$), so it is affine in $k$.

(b) $\delta^*_B(r,s)=\sum_ip_iX_{ir}X_{is}(m_r(X_i,D_i)-m_s(X_i,D_i))$, and every term is
non-negative by hypothesis.

(c) Under M0 with $R=2$, row 1's mean is $b_0+b_1X$ with $b_1=w_1+d_1$ and row 2's is
$w_0+w_1X$. Then $\Cov_{w^1}(b_0+b_1X,X)=b_1\Var_{w^1}(X)$ and likewise for row 2, so
$\delta^*_W=\mu b_1\Var_{w^1}(X)+(1-\mu)w_1\Var_{w^2}(X)$, which is the stated form with
$a=\mu\Var_{w^1}(X)\ge0$ and $b=(1-\mu)\Var_{w^2}(X)\ge0$. Under a district rung, the law
of total covariance under each row weight, conditioning on $D$, gives
\[
\delta^*_W=\sum_s\mu_s\Big(\E_{w^s}\big[\Cov_{w^s}(X,m_s\mid D)\big]+\Cov_{w^s}\big(\E_{w^s}[X\mid D],\E_{w^s}[m_s\mid D]\big)\Big).
\]
Within district $d$, $m_s(\cdot,d)$ is linear with slope $b_{s,d}$, so the conditional
covariance is $b_{s,d}\Var_{w^s}(X\mid D=d)\ge0$ when $b_{s,d}\ge0$. The between-district
term has no sign; under M1 it contains the covariance of the district shifts with
composition. In the six-unit example of the main text the within-district term is
$+0.00064$ and the pooled association $-0.04927$.
\end{proofof}

\emph{Remark~\ref{P-rem:whichdelta}.} Each unit's term
$\sum_rc_{ir}e_{ir}$ has range at most $\sum_r|c_{ir}|$, because each $e_{ir}$ lies in an
interval of length one; by Popoviciu and independence across units, the variance of
$\delta-\delta^*$ is at most $\tfrac14\sum_i(\sum_r|c_{ir}|)^2$.

\subsection{The sign module under the pooled model}

\begin{proofof}{Theorem~\ref{P-thm:R3}}
(a) A value $t$ is in the identified set iff some $\mathbb Q\in\mathcal M(Q(P)+K;A)$ with
$W$-marginal $P$ has $k\in C$ and $\varphi(k)=t$. By (C2), $C$ is evaluated at the true
$P$ for every such $\mathbb Q$, so the condition is $k\in\ThI\cap C$, and the identified
set is $\varphi(\ThI\cap C)$. $\ThI$ is convex and compact (Theorem~\ref{P-thm:sharp},
with $\affdim(A)=R-1$) and $C$ is closed and convex, so $\ThI\cap C$ is convex and
compact or empty, and its image under the affine $\varphi$ is a compact interval or
empty.

(b) ($\supseteq$) For $P'\in\mathcal P(P)$, $Q(P')=Q(P)$ and the design is the same, so
$C(P')=C(P)$ by (C2), and $\ThI(P')\subseteq\ThO(P')=\ThO(P)$. ($\subseteq$) For
$k\in\ThO\cap C$, the construction in the proof of Proposition~\ref{P-prop:sharpreg} gives
$P'\in\mathcal P(P)$ with $k\in\ThI(P')$, and $k\in C(P)=C(P')$. The interval statements
follow from (a) applied to each $P'$ and from Proposition~\ref{P-prop:IequalsO}.

(c) $\varphi(\Theta\cap(C_+\cup C_-))=\varphi(\Theta\cap C_+)\cup\varphi(\Theta\cap C_-)$, and
each piece is a closed interval or empty by (a) or (b). If one piece is empty the
union is the other. A union of two non-empty closed intervals is an interval iff they
intersect, and a point $k_0\in\Theta\cap C_+\cap C_-$ gives $\varphi(k_0)$ in both.

(d) Apply Lemma~\ref{P-lem:M1}(c) with $M_H(X_i)$ in place of $\beta_i$ and $\tau(X_i)$ in
place of $T_i$; the identities use only $T=\sum_sX_s\beta_s$ and the simplex, and
$x^\top M_H(x)=\tau(x)$ holds. This gives
$\delta^{r*}_W=\Var_p(X_r)(\ER^{r*}-D^{r*})$ and
$\delta^{r*}_B=\mu_r(1-\mu_r)(D^{r*}-\NM^{r*})$, where $\NM^{r*}$ is the neighbourhood gap
computed from $\tau$, and $\NM^{r*}=\Cov_p(X_r,\tau)/(\mu_r(1-\mu_r))=\gamma_r\ER^{r*}$. So
both are non-negative iff $\gamma_r\ER^{r*}\le D^{r*}\le\ER^{r*}$, both non-positive iff
$\ER^{r*}\le D^{r*}\le\gamma_r\ER^{r*}$, and the union is
$D^{r*}\in\hull\{\gamma_r\ER^{r*},\ER^{r*}\}$. $D^{r*}$ is affine in $k$ and
$\ER^{r*},\gamma_r$ depend only on the design and $Q$, so this is two rows of an
admissible set. The sign statement follows because $\gamma_r>0$.

(e) $K_0\in\ThI$ by definition and $k_0\in C$ by hypothesis; then apply
Proposition~\ref{P-prop:nesting}.
\end{proofof}

\emph{Remarks~\ref{P-rem:fullreg} and~\ref{P-rem:LCEexcludesCR}.} In
Remark~\ref{P-rem:fullreg}, the full law forces both rows' rates to equal $T$ in every
unit (at each composition $T$ is 0 or 1, so both rates equal $T$), hence both means are
$.5$ and $k=0\notin[.2,.4]$; the regression $\tau=.5$ admits $m_1=.5+k(1-x)$,
$m_0=.5-kx$, valid for $k$ in a range containing $[.2,.4]$; with the target weights
$\E_{w^1}[1-X]=\E_{w^0}[X]=5/12$, so the gap is $\tilde D=5k/6$, which maps $[.2,.4]$
onto $[1/6,1/3]$. In Remark~\ref{P-rem:LCEexcludesCR}, the regression is
$(1-x)^2+(49/55)x^2$, whose cross coefficient in the basis
$((1-x)^2,2x(1-x),x^2)$ is $Q_{12}=0$; with $A=[0,1]$, Proposition~\ref{P-prop:simplex}
gives $|k|\le Q_{12}=0$, so the vertex rows force $k=0$, and
at $k=0$ the computation gives $\ER=0$ and $\tilde D=-261/1925\ne0$, so neither CR branch
contains it.

\emph{A belief can be refutable without the rung.} On $x=1/4,1/2,3/4$ let both rows'
means be $1.5-2x$ and the aggregate equal this function. LCE holds on $[1/4,3/4]$,
$\ER=-2$, and every feasible gap is at least $-1$, so a positive within association is
impossible even in the model-free continuous problem.

\emph{The certificate is only sufficient.} On $[-1,1]^2$ with $g_B=u+.5$ and
$g_W=u-.5$, the branches are the disjoint strips $u\ge.5$ and $u\le-.5$. For the target
$u$ the images are $[.5,1]$ and $[-1,-.5]$; for the target $v$ both images are $[-1,1]$,
although the certificate LP $\{g_B=g_W=0\}$ is infeasible.

\begin{proofof}{Proposition~\ref{P-prop:R2}}
(a) $F\cap N\subseteq F$, so $U_N\le U$. If some $k\in F^*_U\cap N$, then
$U_N\ge\varphi(k)=U$. If $F^*_U\cap N=\emptyset$, then $F\cap N$ is non-empty and compact,
so $\varphi$ attains $U_N$ at some $k'\in F\cap N$, and $k'\notin F^*_U$ gives
$\varphi(k')<U$. The lower end is the same argument for $-\varphi$. For $N=\{Mk\le m\}$,
$F^*_U\cap N=\{Gk\le g,\ \omega^\top k\ge U-\varphi_0,\ Mk\le m\}$, a feasibility LP; the
intercept must be subtracted, or the row $\omega^\top k\ge U$ would reject the optimal face
whenever $\varphi_0\ne0$.

(b) $F^*_U$ is a non-empty compact face, so $\beta^*=\min_{F^*_U}a^\top k$ is attained, and
$F^*_U\cap\{a^\top k\le b\}=\emptyset$ iff $\beta^*>b$; the first claim follows from (a).
Let $v(\beta)=\max\{\varphi(k):k\in F,a^\top k\le\beta\}$ be the perturbation function.
By (a) at $\beta$, $v(\beta)=U$ for $\beta\ge\beta^*$ and $v(\beta)<U$ for $\beta<\beta^*$.
If $b\ge\beta^*$, $v$ is constant on $[b,\infty)$, so $v'_+(b)=0$ and $0$ is an optimal
multiplier. If $b<\beta^*$, concavity gives $v'_+(b)\ge(v(\beta^*)-v(b))/(\beta^*-b)>0$,
so every optimal multiplier, which lies in $[v'_+(b),v'_-(b)]$, is positive (at the left
end of the domain the set is $[v'_+(b),\infty)$). A solver returns one multiplier from
the superdifferential, and when that is $[0,\lambda]$ with $\lambda>0$ it may return a
positive one although the row does not tighten: maximise $x$ on $[0,1]$ with the row
$x\le1$ duplicated, where every split of the unit dual cost is optimal. If the row is
slack at an optimum $k$ of the augmented LP, then near $k$ the augmented set coincides
with $F$, so $k$ is a local, hence global, maximiser of the linear $\varphi$ on $F$, and
$k\in F^*_U\cap N$: the row does not tighten.

(c) If $F\subseteq N$, nothing changes. If a row $a^\top k\le b$ is not implied,
$\max_Fa^\top k>b$, and for the target $a^\top k$, $U_N\le b<U$. For general closed convex
$N$ with $F\not\subseteq N$, take $k_1\in F\setminus N$ and separate it strictly from $N$
by $(a,b)$ with $a^\top k\le b$ on $N$ and $a^\top k_1>b$; for the target $a^\top k$,
$U_N\le b<a^\top k_1\le U$.

(d) With $R=2$, Theorem~\ref{P-thm:R3}(d) gives
$\delta^*_W=\Var_p(X)(\ER^*-\tilde D)$ and $\delta^*_B=\mu(1-\mu)(\tilde D-\NM^*)$, and
$\tilde D$ is affine in the scalar unknown with slope proportional to $-(1-\gamma)$, so
non-constant iff $\gamma<1$. Each sign row is a half-line in $\tilde D$, and the feasible
set maps bijectively onto $[L_J,U_J]$; cutting the set therefore means cutting the
interval, so ``cuts the set only'' is impossible for $\tilde D$ and for any target
non-constant in the unknown. Intersecting with the half-lines gives the closed forms.

(e) This is (a) applied to the stack and to each sub-stack; the optimal face changes as
rows are added, so contributions depend on order. Example: on $[0,1]^2$ with target $x$
and rows $x+y\le.8$, $x\le.5$, adding $x+y\le.8$ last gives zero gain for $x$, yet it
removes $(.5,1)$ and lowers the upper end of $x+y$ from 1.5 to .8.
\end{proofof}

\begin{proofof}{Proposition~\ref{P-prop:G2}}
EGGHM's identities, $D_{\ER}-D=\delta_W/\Var_p(X)$ and $D_{\NM}-D=-\delta_B/(\mu(1-\mu))$,
are accounting identities for any mean function (Lemma~\ref{P-lem:M1}), so they hold for
the model-implied quantities. With $\Var_p(X)>0$ and $0<\mu<1$, $\delta^*_W\ge0\iff\tilde
D\le\ER^*$ and $\delta^*_B\ge0\iff\tilde D\ge\NM^*$. Under LCE, $\tilde D$ is affine and
non-constant in the unknown slope when $\gamma<1$, so the JKST set maps onto $[L_J,U_J]$.
Intersecting with the half-lines gives the four cases, and CR is the union of the two
same-sign cases, that is, $\tilde D$ between $\NM^*$ and $\ER^*$.
\end{proofof}

\begin{proofof}{Proposition~\ref{P-prop:sandwich}}
$\tau(x)=x^\top Hx$, so $\nabla\tau=(H+H^\top)x=2Qx$ and
$\partial_{e_r-e_s}\tau=2q_{rs}$. Also $x^\top H=(H^\top x)^\top$ with $H^\top=Q-K$. So
$O_{rs}=(Qx)_r-(Qx)_s+(Kx)_r-(Kx)_s=q_{rs}+g_{rs}$ and $L_{rs}=q_{rs}-g_{rs}$, which is
(b); adding gives (a); and (c) follows from (b).
\end{proofof}

\begin{proofof}{Proposition~\ref{P-prop:rank}}
(a) $g_{rs}(x;k)=\sum_t(K_{rt}-K_{st})x_t$. As $x$ ranges over a full-dimensional subset of
the simplex, the functionals reached are $K_{rs}$ and $K_{rt}-K_{st}$ for $t\ne r,s$:
$R-1$ independent directions. (b) For a star $\{(1,s)\}$, $K_{1s}$ is reached directly,
and then $K_{st}=K_{1t}-(K_{1t}-K_{st})$. (c) follows from (b) and
Proposition~\ref{P-prop:sandwich}(c), which bounds each reached functional by an
identified quantity. Rank is necessary, not sufficient, for tightening. For general
spanning trees the source argument (induction on the leaves of the tree) is only
outlined, and the claim has been checked numerically for chains only.
\end{proofof}

\begin{proofof}{Proposition~\ref{P-prop:PN}}
(a) $w_1+d_1x=(1-x)w_1+x(w_1+d_1)$, a convex combination of the two slopes for
$x\in[0,1]$. (b) From the proof of Proposition~\ref{P-prop:Krows}(c),
$\delta^*_W=a(w_1+d_1)+bw_1=(a+b)(w_1+d_1x^*)$ with $x^*=a/(a+b)$. So if
$w_1+d_1x\ge0$ on the hull and $x^*$ lies in the hull, $\delta^*_W\ge0$.
\end{proofof}

\emph{On Conjecture~\ref{P-conj:C1}.} Writing
$\Var_w(X)=\tfrac12\sum_{i,j}w_iw_j(X_i-X_j)^2$ expresses $x^*$ as a weighted mean of
pair ratios, but individual ratios can leave the hull, so a proof needs a cancellation
argument that has not been found.

\subsection{Monotone nonparametric contextual effects}

\begin{proofof}{Lemma~\ref{P-lem:N21}}
By the tower property with A3 and conditional Jensen for $\Bset$,
$m(X)=\E[\beta\mid X]\in\Bset$ a.s. And $m(X)^\top X=\E[\beta^\top X\mid X]=\E[T\mid
X]=\tau(X)$ by A1. On a finite support every point has positive mass, so these hold at
every support point.
\end{proofof}

\begin{proofof}{Proposition~\ref{P-prop:N2}}
(a) Every N2($\sigma$) mean satisfies Lemma~\ref{P-lem:N21} at each $x_j$ and the order
rows on the support, and the target depends on $m$ only at the support points.
(b) Extend the whole matrix $m$ as a step function, constant between support points.
The extension is Borel, each row stays in the simplex, each designated entry stays
monotone, and its values at the support points are unchanged. On a strict subset of
the support this fails: with true support $\{1/4,1/2,3/4\}$, row-1 means $(.4,.4,.4)$,
row-2 means $(8/15,.7,.8)$ and $\tau=(.5,.55,.5)$, a grid keeping only the extremes
accepts both rows equal to $.5$ at both extremes, and any monotone extension then
assigns $.5$ to both rows at the middle, which cannot give $.55$.
(c) For $K\in\ThO$, $M_{Q+K}(x)\ge0$ on $A\supseteq\supp P_X$ and its rows sum to one, so it
lies in $\Bset$; and $x^\top M_{Q+K}(x)=x^\top Qx=\tau(x)$, so it lies in $S(x,\tau(x))$ at
each atom. A linear function with a slope of sign $\sigma$ is $\sigma$-monotone.
(d) $\delta^*_W=\mu\Cov_{w^1}(m_1(X),X)+(1-\mu)\Cov_{w^2}(m_2(X),X)$, and both
covariances are non-negative by Chebyshev's association inequality applied under $w^1$
and under $w^2$.
(e) With the chart $(x_1,x_3)$, so that $x_2$ adjusts, $m_{rc}(x)=x_3$ is constant, hence
weakly non-decreasing, in $x_1$; under a symmetric Dirichlet design
$\E[X_3\mid X_1]=(1-X_1)/2$, and tilting by $X_1$ or $1-X_1$ preserves that conditional
mean, so the covariance with $X_1$ is negative under both W1 weights. A globally
increasing univariate function of $X_1$ still obeys the covariance inequality.
\end{proofof}

\emph{On Conjecture~\ref{P-conj:N2}.} In the size example of the main text, at the
first size $T\in\{0,1\}$, so both rows' rates equal $T$ and their conditional means are
$(.5,.5)$; at the second size $T=1/2$ and a selection $(1,0)$ is in $S(1/2,1/2)$. The
polytope LP, conditioning on $X$ only, therefore admits $(.75,.25)$, which the model
excludes. \emph{Sufficiency on a finite support} (a sketch of the argument that exists):
if $m(x)$ lies in the selection expectation of $S(x,T)$ given each positive-mass design
atom $(x,n)$, take for each atom a measurable selection whose conditional mean is $m(x)$
and glue the selections across the finitely many atoms; the accounting identity holds
a.s.\ and the shape restrictions hold because they concern the common $m$. Necessity in
general, through the selection construction of Theorem~\ref{P-thm:sharp} applied to each
admissible $m$, is not written.

\begin{proofof}{Proposition~\ref{P-prop:N23}}
(a) Differentiate $\tau_{m,c}(z)=zm_{1c}(z)+(1-z)m_{2c}(z)$ by the product rule; (b) is
a rearrangement. (c) If $\tau'\ge0$: $D\ge0$ and $\delta_W\ge0$ give
$0\le D=\tau'-\delta_W\le\tau'$; $D\le0$ and $\delta_W\le0$ give
$D=\tau'-\delta_W\ge\tau'\ge0\ge D$, so $D=0$; conversely $D\in[0,\tau']$ gives $D\ge0$ and
$\delta_W=\tau'-D\ge0$. The case $\tau'\le0$ is symmetric. (d) $x,1-x\ge0$, so
non-negative derivatives give $\delta_W(x)\ge0$. Both rows equal to $.5+h$ with
$h=(x-\tfrac14)(x-\tfrac12)(x-\tfrac34)$ are constant on $\{1/4,1/2,3/4\}$ but have
derivative $-1/16$ at $1/2$. With $X$ uniform on $[1/5,2/5]$, $m_1=x-x^2/2$ and
$m_0=1-x^2/2$ are probabilities, $xm_1'+(1-x)m_0'=0$ everywhere, and
$\delta^*_W=-1/47250$, so the local sign everywhere does not imply the aggregate sign.
For the reverse, complement both functions, shrink them slightly into the interior and
subtract a small positive multiple of $x$ from each: the aggregate within association
stays positive by continuity and the local one becomes negative.

\emph{Not identified on a finite support.} Adding $h(x)=\prod_j(x-x_j)$ to both rows
leaves the law on the support unchanged and moves every nodal derivative. With support
$\{1/4,1/2,3/4\}$, $m_1=.7$, $m_0=.3$ and $h$ as above, both versions are valid on $[0,1]$
and induce the same law, but $\tau'(1/2)$ is $.4$ in one and $.3375$ in the other.
\end{proofof}

\begin{proofof}{Lemma~\ref{P-lem:bins}}
(i) $\bar x_b\bar m_{1b}+(1-\bar x_b)\bar m_{2b}=\E_b[Xm_1+(1-X)m_2]=\E_b\tau=\bar\tau_b$ by
Lemma~\ref{P-lem:N21}, and each $\bar m_{rb}$ is a weighted average of simplex points.
(ii) $\E_p[Xm_1]=\sum_bP(b)\E_b[Xm_1]=\sum_bP(b)\bar x_b\bar m_{1b}$, and likewise for row 2.
(iii) $\bar m_{1b}$ is an average of $m_1$ over bin $b$ under weights proportional to
$pX$, so it lies in $[\min_bm_1,\max_bm_1]$; ordered non-overlapping bins and monotone
$m_1$ give $\max_bm_1\le\min_{b+1}m_1$. Row 2 is the same with weights proportional to
$p(1-X)$. A group absent from a bin is given any value that keeps the bin rates
monotone. For the two-point example, equal bin weights at $x=.25,.75$ with $m_1=x$ and
$m_2=0$ give $\E_bT=\E[X^2]=.3125$; the ordinary bin-average rates evaluated at
$\E X=.5$ give $.25$, while the person-weighted rate is $.625$ and reproduces $.3125$.
\end{proofof}

\emph{Degree selection.} In a balanced two-level design with true slope zero, choosing
the slope by leave-one-out error selects it when $Z^2>2$, where $Z$ is the slope's
$t$-statistic; the nominal 95\% interval then covers with conditional probability
$1-.05/P(|Z|>\sqrt2)\approx0.68$.
